% Source: Solutions/Dhrubojoti-PS4-Q1-4.pdf, pp. 4-9 (typed). Restated problem statement omitted.
\submittedby{Dhrubojoti}

\textbf{Recall.}
A map $f\colon X\to Y$ is continuous if and only if the preimage of every closed subset of $Y$
is closed in $X$. Also, if $A$ is a closed subset of $X$, then every relatively closed subset
of $A$ is closed in $X$.

\paragraph{(i) The collection $\{A_\alpha\}$ is finite and each $A_\alpha$ is closed.}
\begin{proof}
Suppose
\[ X=A_1\cup A_2\cup\dots\cup A_n, \]
where each $A_i$ is closed in $X$. We are given that the restriction
\[ f|_{A_i}\colon A_i\to Y \]
is continuous for every $i=1,2,\dots,n$. We need to prove that $f\colon X\to Y$ is continuous.

Let $C$ be an arbitrary closed subset of $Y$. Consider
\[ f^{-1}(C)=\{x\in X : f(x)\in C\}. \]
Since $X=\bigcup_{i=1}^n A_i$, we can write
\[ f^{-1}(C)=\bigcup_{i=1}^n\bigl(f^{-1}(C)\cap A_i\bigr). \]
Now observe that
\[ f^{-1}(C)\cap A_i=(f|_{A_i})^{-1}(C). \]
Since $f|_{A_i}$ is continuous and $C$ is closed in $Y$, $(f|_{A_i})^{-1}(C)$ is closed in
$A_i$. But $A_i$ itself is closed in $X$. Therefore, a relatively closed subset of $A_i$ is
also closed in $X$. Thus, $f^{-1}(C)\cap A_i$ is closed in $X$ for every $i$. Consequently,
\[ f^{-1}(C)=\bigcup_{i=1}^n\bigl(f^{-1}(C)\cap A_i\bigr) \]
is a finite union of closed subsets of $X$. Hence, $f^{-1}(C)$ is closed in $X$. Since $C$ was
an arbitrary closed subset of $Y$, we conclude that
\[ \boxed{f\colon X\to Y\text{ is continuous.}} \qedhere \]
\end{proof}

\paragraph{(ii) Find a countable collection of closed sets covering $X$ for which $f$ is not continuous.}

\textbf{Example.} Take
\[ X=\R,\qquad Y=\R. \]
Define
\[ A_0=\{0\} \]
and, for $n\in\N$,
\[ A_n=\left[-n,-\frac1n\right]\cup\left[\frac1n,n\right]. \]

\textbf{Step 1. Show that each $A_n$ is closed.}
Clearly, $A_0=\{0\}$ is closed in $\R$. Also, each $A_n$ is a finite union of closed intervals.
Hence, every $A_n$ is closed in $\R$.

\textbf{Step 2. Show that the collection is countable and covers $\R$.}
The collection
\[ \{A_0,A_1,A_2,\dots\} \]
is countable. Moreover, every nonzero real number belongs to some $A_n$. Indeed, given
$x\neq0$, choose $n$ sufficiently large such that
\[ n\ge|x|\quad\text{and}\quad\frac1n\le|x|. \]
Then $x\in A_n$. Since $0\in A_0$, we obtain
\[ \boxed{\R=\bigcup_{n=0}^\infty A_n.} \]

\textbf{Step 3. Define the map $f$.}
Define $f\colon\R\to\R$ by
\[ f(x)=\begin{cases} 1, & x=0,\\ 0, & x\neq0.\end{cases} \]
We now show that the restriction of $f$ to every $A_n$ is continuous.
\begin{itemize}
  \item On $A_0=\{0\}$, the restriction is the constant function $f(0)=1$, which is continuous.
  \item On every $A_n$, $n\ge1$, we have $0\notin A_n$. Therefore, $f|_{A_n}$ is the constant
        zero function, which is continuous.
\end{itemize}
Thus, every restriction $f|_{A_n}\colon A_n\to\R$ is continuous.

\textbf{Step 4. Show that $f$ is not continuous.}
Consider the open set
\[ U=\left(\frac12,\frac32\right)\subseteq\R. \]
We have
\[ f^{-1}(U)=\{0\}. \]
But $\{0\}$ is not open in $\R$. Therefore, $f$ is not continuous.

Hence, we have constructed a countable closed cover of $\R$ such that every restriction is
continuous, but the original map is not continuous. Thus,
\[ \boxed{\text{The conclusion of (i) need not hold for a countable cover.}} \]
\hfill$\square$

\paragraph{(iii) Prove the result when $\{A_\alpha\}$ is locally finite and each $A_\alpha$ is closed.}
\begin{proof}
We are given that:
\begin{enumerate}
  \item $X=\bigcup_\alpha A_\alpha$.
  \item Every $A_\alpha$ is closed in $X$.
  \item The family $\{A_\alpha\}$ is locally finite.
  \item Every restriction $f|_{A_\alpha}\colon A_\alpha\to Y$ is continuous.
\end{enumerate}
We need to prove that $f\colon X\to Y$ is continuous. We will use the following important
result.

\begin{lemma}
The union of a locally finite family of closed subsets of a topological space is closed.
\end{lemma}

\emph{Proof of the lemma.}
Let $\{F_\alpha\}$ be a locally finite family of closed subsets of $X$. We want to show that
\[ F=\bigcup_\alpha F_\alpha \]
is closed in $X$. Take an arbitrary point $x\in X\setminus F$. Since the family is locally
finite, there exists an open neighbourhood $U$ of $x$ that intersects only finitely many members
of the family. Suppose these members are
\[ F_{\alpha_1},F_{\alpha_2},\dots,F_{\alpha_n}. \]
Now, $x\notin F$, so in particular,
\[ x\notin F_{\alpha_i},\qquad i=1,\dots,n. \]
Since each $F_{\alpha_i}$ is closed, its complement is open. Thus, $X\setminus F_{\alpha_i}$ is
an open neighbourhood of $x$. Consider
\[ V=U\cap\bigcap_{i=1}^n(X\setminus F_{\alpha_i}). \]
This is an open neighbourhood of $x$, being a finite intersection of open sets. We claim that
\[ V\cap F=\emptyset. \]
Indeed:
\begin{itemize}
  \item For the finitely many sets $F_{\alpha_1},\dots,F_{\alpha_n}$, the set $V$ does not
        intersect any of them.
  \item For every other index $\alpha$, $U$, and hence $V$, does not intersect $F_\alpha$, by
        the choice of $U$.
\end{itemize}
Therefore, $V$ does not intersect their union $F$. So every point in $X\setminus F$ has an open
neighbourhood contained in $X\setminus F$. Hence, $X\setminus F$ is open, which means
\[ \bigcup_\alpha F_\alpha\text{ is closed in }X. \]
This proves the lemma.

Now we prove the required result. Let $C$ be an arbitrary closed subset of $Y$. Consider
$f^{-1}(C)$. Since $X=\bigcup_\alpha A_\alpha$, we have
\begin{align*}
  f^{-1}(C) &= \bigcup_\alpha\bigl(f^{-1}(C)\cap A_\alpha\bigr)\\
            &= \bigcup_\alpha(f|_{A_\alpha})^{-1}(C).
\end{align*}
Now, since every restriction $f|_{A_\alpha}$ is continuous and $C$ is closed in $Y$,
$(f|_{A_\alpha})^{-1}(C)$ is closed in $A_\alpha$. Since $A_\alpha$ is closed in $X$, it follows
that
\[ F_\alpha:=(f|_{A_\alpha})^{-1}(C) \]
is closed in $X$. Furthermore, since $F_\alpha\subseteq A_\alpha$ and the family $\{A_\alpha\}$
is locally finite, the family $\{F_\alpha\}$ is also locally finite. Therefore, by our lemma,
the union
\[ f^{-1}(C)=\bigcup_\alpha F_\alpha \]
is closed in $X$. Since $C$ was an arbitrary closed subset of $Y$, the preimage of every closed
subset of $Y$ is closed in $X$. Consequently,
\[ \boxed{f\colon X\to Y\text{ is continuous.}} \qedhere \]
\end{proof}
